% Derived from leanblueprint project templates; modified for AbsorptionCutoff.
% In this file you should put macros to be used only by
% the printed version. Of course they should have a corresponding
% version in macros/web.tex.
% Typically the printed version could have more fancy decorations.
% This should be a very short file.
%
% This file starts with dummy macros that ensure the pdf
% compiler will ignore macros provided by plasTeX that make
% sense only for the web version, such as dependency graph
% macros.


% Make the leanblueprint status flags visible in the printed version, so the
% PDF is a self-contained snapshot of formalization progress (the web version
% shows these via dependency-graph coloring instead). Print-only: keep the
% matching macros in macros/web.tex as-is.
\usepackage{xcolor}
\usepackage{fontspec}
% The web blueprint exposes every declaration attached by \lean{...}, and the
% generated lean_decls file audits those names.  Printing the raw identifiers
% in the narrow page margin makes long nodes unreadable, so the PDF suppresses
% the names and retains the inline [Lean] status badge below.
\newcommand{\lean}[1]{}
\newcommand{\discussion}[1]{}
% Inline status badges (no argument): rendered where they appear in source.
\newcommand{\bpbadge}[2]{{\footnotesize\textcolor{#1}{[#2]}}\ }
\newcommand{\leanok}{\bpbadge{green!55!black}{Lean}}
\newcommand{\mathlibok}{\bpbadge{blue!70!black}{Mathlib}}
\newcommand{\notready}{\bpbadge{orange!85!black}{todo}}
% Make sure that arguments of \uses and \proves are real labels, by using invisible refs:
% latex prints a warning if the label is not defined, but nothing is shown in the pdf file.
% It uses LaTeX3 programming, this is why we use the expl3 package.
\ExplSyntaxOn
\NewDocumentCommand{\uses}{m}
 {\clist_map_inline:nn{#1}{\vphantom{\ref{##1}}}%
  \ignorespaces}
\NewDocumentCommand{\proves}{m}
 {\clist_map_inline:nn{#1}{\vphantom{\ref{##1}}}%
  \ignorespaces}
\ExplSyntaxOff
